\chapter{Genomic organization, regulation and development}
\ChapterGoals{%
\item describe gene expression as a physical mechanism from DNA to RNA to protein;
\item distinguish cis and trans regulation, and chromosome context from nearest-gene lookup;
\item nondimensionalise a two-regulator model and derive its fixed points and their stability;
\item explain how history-dependent states arise without new parameters;
\item separate recorded lineage from inferred order, and within-cell regulation from population selection.}
{Chapter~7 (rate equations); Appendix~\gref{app:biology} (central dogma).}
Take two cells with the same inherited DNA. Why can one produce a regulator while
the other does not? The DNA sequence is not the entire initial condition.
Factor concentrations, chromatin, signals and history change the molecular
processes that read it. Chapter 29 optimized explicit mathematical objects;
this chapter asks what biological objects a genomic model would have to represent.

We separate locus-linked sequence, diffusible regulators, molecular state,
observed time, ancestry and population change. Then we build a minimal regulatory
model, derive its alternative stable states, and deliberately break several
tempting interpretations. The model is a transparent teaching artifact, not a
reconstruction of an organism or a replication of a synthetic circuit experiment.

\section{Start with the physical expression mechanism}
A protein-coding locus supplies a transcription template and regulatory sequence.
Transcription machinery makes RNA; processing and transport may intervene;
ribosomes translate suitable messenger RNA into polypeptides. RNA genes need not
produce proteins. DNA recognition, factor binding, transcription initiation,
elongation, RNA abundance and protein abundance are separate events or variables.
\Plate{01-expression}{An original eukaryotic expression schematic separates
regulatory binding, transcription, processing/export and translation.
DNA is not translated directly; the drawing omits detailed molecular structures
and does not depict a universal prokaryotic organization.}{dnad30:expression}

A minimal observation-aware model is
\begin{align}
 \dot m&=q(G,R,E)-\gamma_m m,\\
 \dot p&=\beta m-\gamma_p p.
\end{align}
Here $m$ and $p$ are RNA and protein concentrations; $q$ has concentration/time
units, $\beta,\gamma_m,\gamma_p$ have inverse-time units, and $G,R,E$ denote
sequence, cellular/regulatory state and environment. The RNA and protein
concentration units may have different species reference scales, in which case
$\beta$ includes their conversion. For a fixed nonnegative transcription rate,
$m^\star=q/\gamma_m$ and $p^\star=\beta q/(\gamma_m\gamma_p)$ when both loss rates
are positive. Protein abundance alone cannot identify all four parameters.
A large protein pool can persist after transcription changes.

This is not a complete expression model: maturation, bursting, transport,
translation regulation, cell growth and division may matter. Expression
must name its measured layer. Fluorescence from a reporter adds yet another
observation transfer function. Chapter 10's readout discipline still applies.

\section{Cis and trans are relational distinctions}
A \Term{cis-regulatory element} is linked to the locus it influences.
A \Term{trans regulator}, such as a diffusible transcription factor, can act on
multiple compatible loci. Promoters support transcription initiation; enhancers
can affect expression through context-dependent communication with promoters.
Neither term means a universal software opcode.

Consider two alleles exposed to one shared factor concentration $x$.
For one stipulated independent binding site, let free-site and bound-site weights
be $1$ and $x/K$, with dissociation scale $K>0$. Normalization gives
\[
 Z=1+x/K,\qquad \pi_{\rm bound}=\frac{x/K}{Z}=\frac{x}{K+x}.
\]
This equilibrium occupancy assumes an available site, constant free factor and
a specified binding model. It is not the probability that an arbitrary sequence
motif is transcriptionally active.

For illustration only, assign allele-specific maximum rates $c_i$ and
accessibility multipliers $q_i\in[0,1]$:
\[
 u_i=c_iq_i\frac{x}{K+x},\qquad i\in\{A,B\}.
\]
Here $c_i$ has production-rate units; $q_i$ and occupancy are dimensionless.
The multiplicative accessibility rule is stipulated, not a general biochemical
law. Changing $c_A$ changes only A in this model; changing shared $x$ changes both.
\Listing{allele_rates}
\begin{center}\input{\Generated/alleles.tex}\end{center}
With $K=1$, equal cis rates and open sites, $x=1$ gives $(.5,.5)$.
Halving only A's cis rate gives $(.25,.5)$; increasing shared factor to three gives
$(.75,.75)$. Closing only A's site gives $(0,.5)$.
Allelic imbalance can suggest a locus-linked difference, but mapping,
copy-number, imprinting and measurement biases require controls. Equal responses
do not prove identical cis mechanisms.

\section{Chromosome context is not nearest-gene lookup}
DNA is packaged with proteins. Local nucleosome organization, accessibility,
factor availability and chromosome contacts contribute different information.
Contact is not occupancy; accessibility is not productive transcription.
An enhancer's nearest linear gene is not its target by definition.
\Plate{02-context}{A distal element and a promoter can be close in a folded
chromosome despite an intervening locus. A dashed contact link is not a covalent
bond or an asserted causal regulatory effect. Two allele panels isolate a
local cis change from a shared diffusible-factor change.}{dnad30:context}

Fulco et al. tested candidate enhancer--gene connections using CRISPR
perturbations and formulated an activity-by-contact predictor
\cite{Fulco2019}. This provides a concrete example of connecting a regulatory
prediction to perturbation evidence. It does not make accessibility times contact
a universal exact transcription law. Our two-allele calculation does not implement
that predictor or its experiments.

For a mechanistic hypothesis, identify the intervention, locus, cell context,
measurement window and alternative explanation. Perturbing an element can affect
multiple local processes; removing a factor can have indirect consequences.
Independent occupancy, transcription and phenotype measurements address
different links in the proposed chain. A correlation graph alone does not
identify those links.

Maizels and Briscoe's 2026 perspective distinguishes statistical GRN descriptions
from mechanistic explanations and argues for biologically constrained models
tested through richer perturbations \cite{MaizelsBriscoe2026}.
This is a research agenda, not evidence that an arbitrary learned latent state
has acquired causal biological meaning. Our exact toggle calculation illustrates
one declared model; it does not infer a regulatory network from single-cell data.

\section{Nondimensionalize a two-regulator model}
Suppose regulator A represses B production and B represses A production.
Choose a phenomenological Hill-$2$ production term and a common first-order loss:
\[
 \frac{dA}{dt}=\frac{v_A s}{1+(B/K)^2}-\delta A,\qquad
 \frac{dB}{dt}=\frac{v_B}{1+(A/K)^2}-\delta B.
\]
$A,B,K$ are concentrations, $v_A,v_B$ are concentration/time, $s\ge0$ is a
dimensionless signal multiplier, and $\delta>0$ is inverse time.
Set $a=A/K$, $b=B/K$, $\tau=\delta t$ and
$\alpha_i=v_i/(\delta K)$:
\[
 \frac{da}{d\tau}=\frac{\alpha_A s}{1+b^2}-a,\qquad
 \frac{db}{d\tau}=\frac{\alpha_B}{1+a^2}-b.
\]
All reference states and times below use these dimensionless coordinates.
The genotype record stores abstract cis-dependent strengths, not nucleotide
sequences. A measured concentration cannot be inserted without specifying $K$.

The exponent is an effective response-shape choice. It does not establish that
exactly two proteins bind cooperatively. Binding equilibrium, RNA intermediates,
resource competition, spatial structure and molecular noise are omitted.
Gardner, Cantor and Collins constructed a mutually inhibitory synthetic toggle
in \emph{E. coli} \cite{Gardner2000}; our parameter values, equations and
trajectories are not a reproduction of their device.

\section{Derive fixed points rather than naming every plateau memory}
Choose $s=1$ and $\alpha_A=\alpha_B=\alpha$. At a fixed point,
\[
 a(1+b^2)=\alpha,\qquad b(1+a^2)=\alpha.
\]
Subtract to obtain
\[
 (a-b)(1-ab)=0.
\]
There are two cases. A symmetric point satisfies $a=b=c$ and
$c^3+c=\alpha$, with one nonnegative root because $3c^2+1>0$.
An asymmetric point satisfies $ab=1$; substitution gives $a+b=\alpha$.
Thus
\[
 a,b=\frac{\alpha\pm\sqrt{\alpha^2-4}}2
\]
in opposite order, provided $\alpha>2$. At $\alpha=2$ these coincide with
$(1,1)$; they are not three different points. For $\alpha<2$ no real
asymmetric positive pair exists.
\par\medskip\noindent\begin{minipage}{\linewidth}
\Listing{fixed_points}
\end{minipage}\par\medskip
Bisection uses the symmetric polynomial's monotonicity; the asymmetric roots
are algebraic. The reciprocal small root avoids subtracting nearby large numbers.
The teaching implementation restricts this fixed-point calculation to
$0\le\alpha\le100$.

\section{Stability is a property of dynamics, not a picture}
The Jacobian in dimensionless time is
\[
 J=\begin{pmatrix}
 -1 & -2\alpha b/(1+b^2)^2\\
 -2\alpha a/(1+a^2)^2 & -1
 \end{pmatrix}.
\]
At the symmetric point define $u=2c^2/(1+c^2)$.
Eigenvectors $(1,1)$ and $(1,-1)$ have eigenvalues $-1-u$ and $-1+u$.
Common-mode perturbations decay; antisymmetric perturbations grow when $c>1$,
equivalently $\alpha>2$. At the threshold one eigenvalue is zero and a linear
test alone does not resolve the nonlinear boundary behavior.

At an asymmetric point, $ab=1$ and $a+b=\alpha$ give
\[
 \left(\frac{\partial\dot a}{\partial b}\right)
 \left(\frac{\partial\dot b}{\partial a}\right)=\frac4{\alpha^2},
 \qquad \lambda_\pm=-1\pm\frac2\alpha.
\]
Both eigenvalues are negative for $\alpha>2$, establishing local asymptotic
stability. At $\alpha=3$ the values are $-5/3,-1/3$.
\begin{center}\input{\Generated/fixed-points.tex}\end{center}
\Plate{03-network}{The mutual-repression circuit motivates the stipulated
equations; its phase-plane nullclines meet at two locally stable asymmetric
points and one symmetric saddle for $\alpha=3$. Open versus filled symbols
encode stability without relying on color.}{dnad30:network}
Local stability is not a proof of every global trajectory or a measured cell fate.
Memory here means history-dependent persistence in a declared deterministic
dynamical system. It is not evidence of genome editing or a complete epigenetic
mechanism.

\section{Implement positive updates without clipping failure away}
Euler uses old-state production on both sides:
\[
 a_{j+1}=(1-h)a_j+h\frac{\alpha_A s}{1+b_j^2},\qquad
 b_{j+1}=(1-h)b_j+h\frac{\alpha_B}{1+a_j^2}.
\]
For $0<h\le1$, these are nonnegative sums. If $M_A=\alpha_A s$,
then $a_{j+1}\le\max(a_j,M_A)$; similarly
$b_{j+1}\le\max(b_j,\alpha_B)$. This proves an invariant bounded nonnegative
rectangle, not accuracy at $h=1$.
\Listing{euler_step}
Updating A first and using that new A for B is a different algorithm.
A test evaluates both coordinates against independent rational arithmetic and
rejects the sequential-update mutant. At $a=b=10,h=2$, an unprotected Euler
step is negative. Replacing it by zero would conceal an invalid numerical scheme;
the reference rejects the step instead.

\begin{center}\input{\Generated/refinement.tex}\end{center}
At $\tau=10$, successive halved steps converge toward an independently implemented
dense RK4 oracle whose RHS does not call the reference Hill function.
The error ratios lie between 1.8 and 2.2 for these refinements, consistent with
first-order Euler convergence. Positivity, convergence and biological adequacy
remain three separate questions.

\section{Same parameters, different history}
With $\alpha=3$, $h=.01$, initial $(.11,.10)$ and swapped $(.10,.11)$
approach opposite stable branches:
\begin{center}\input{\Generated/terminal.tex}\end{center}
The exactly symmetric initial state stays exactly symmetric in this deterministic
implementation, even though its limiting point is a saddle. It never spontaneously
chooses a branch. Real finite-copy noise is not represented by that exact symmetry.

A production-off intervention sets A's production term to zero; it does not
instantly set A concentration to zero. Its continuous decay is
$a(\tau)=a_0e^{-\tau}$, while Euler gives exactly $a_j=a_0(1-h)^j$.
The tests check the discrete formula independently.
B's response is a causal consequence \emph{inside the stipulated model}, not
proof of a biological A-to-B linkage. A concentration clamp, deletion, inhibition
of transcription and protein degradation are different interventions.

\section{Development, ancestry and inferred order}
Development changes a multidimensional cell state through signaling, regulation,
growth, mechanics and division. A state-space branch is not automatically a
cell-division tree. Observed chronological time, \Term{pseudotime} inferred from
snapshots, and \Term{lineage} ancestry answer different questions.
\Plate{04-development}{The same three supplied cells have observed parent--child
edges but a different pseudotime ordering. The molecular partition below shows
why identical DNA inheritance does not require equal daughter molecular state.
These are teaching records, not reconstructed biological data.}{dnad30:development}
Our records are root $(t=0,z=.2)$, left $(t=1,z=.9)$ and right $(t=1,z=.1)$.
The ancestry edges are root-to-left and root-to-right, yet pseudotime sorting gives
right, root, left. Linking consecutive sorted cells would invent right-to-root
ancestry and reverse observed time.
\Listing{lineage_edges}
The validator uses unique IDs, known parents and strictly earlier parent observation
times. This deliberately restricted tree/forest schema excludes equal-time ancestry
records; it is not a universal lineage-data format or a pseudotime estimator.
Weinreb et al. combine expressed barcodes and transcriptomes to link state and
fate \cite{Weinreb2020}. That experimental linkage is distinct from sorting
snapshot expression.

\section{Inheritance has more than one carrier}
Copying inherited DNA and partitioning cytoplasmic molecules are different
operations. In an ideal symmetric division with independent placement of each of
$N$ identical molecules, daughter-A count $X$ follows
\[
 P(X=k)={N\choose k}2^{-N},\quad
 \mathbb E[X]=N/2,\quad \operatorname{Var}(X)=N/4.
\]
Daughter B receives $N-X$, not an independent binomial draw.
Their count covariance is $-N/4$.
If each daughter has half the original volume, their expected concentrations equal
the parent's, while fluctuations remain. Active segregation, unequal volume and
regulatory re-establishment require a different model.
\begin{center}\input{\Generated/partition.tex}\end{center}
For $N=3$ the exact probabilities are $1/8,3/8,3/8,1/8$.
Enumerating all eight placement bitstrings supplies an independent oracle.
The local seeded simulator conserves molecule count and does not mutate the global
random stream. It is not a model of chromosome segregation.

\section{Population selection is not within-cell adaptation}
Natural selection operates on heritable variation with differential reproductive
success. Let $p$ be type-1 frequency before a haploid reproduction step,
and $w_1,w_0\ge0$ be expected reproductive weights in a specified environment.
With positive mean weight $\bar w=pw_1+(1-p)w_0$,
\[
 p_{\rm sel}=\frac{pw_1}{\bar w}.
\]
The difference from the old frequency is
$p_{\rm sel}-p=p(1-p)(w_1-w_0)/\bar w$.
A common positive weight scale changes no frequency.
With type-1-to-0 mutation probability $\mu$ and reverse probability $\nu$,
\[
 p'=(1-\mu)p_{\rm sel}+\nu(1-p_{\rm sel}).
\]
\Listing{selection_mutation}
Selection followed by mutation generally differs from mutation followed by
selection. The tests explicitly expose this order dependence.
With $p_0=.25,w_1=2,w_0=1,\mu=\nu=0$, odds multiply by two each generation:
\[
 \frac{p_t}{1-p_t}=2^t\frac{p_0}{1-p_0},\qquad
 p_t=\frac{2^t}{2^t+3}.
\]
\Plate{05-evolution}{One type-1 and three type-0 parents generate two and three
descendants under stipulated reproductive weights. Population frequency changes
from $1/4$ to $2/5$. Regulatory feedback in one cell is a different process.}{dnad30:evolution}
\begin{center}\input{\Generated/selection.tex}\end{center}
This deterministic infinite-population expectation excludes drift, linkage,
recombination, changing environments and density dependence. Its weights are
invented, not measured organismal fitness. Gradient descent, transcriptional
adaptation and selecting neural checkpoints do not thereby become biological
natural selection.

\section{Carry biology into computation through explicit maps}
Separate inherited sequence $G$, cellular state $R_t$, environment $E_t$ and
observations $Y_t$:
\[
 R_{t+1}=F(G,R_t,E_t,\xi_t),\qquad Y_t=H(G,R_t,E_t,\eta_t).
\]
Noise in dynamics and measurement has different roles. A deterministic
$G\mapsto Y$ lookup cannot generally express context and history.
\Plate{06-abstraction}{A biological regulatory carrier, a measured variable and a
computational operator are different objects. Mapping concentration to hidden
state preserves a chosen dependency while omitting chemistry, partition and
measurement; an analogy must expose its omitted mechanisms.}{dnad30:abstraction}

If a model calls a gate a promoter, ask what formal operator changed relative to
an established gated recurrence. If it calls hidden state epigenetic memory,
identify persistence, reset and inheritance semantics rather than assuming molecular
fidelity. Testing such computational hypotheses belongs to software research; this
chapter's task is the biological and molecular boundary. A useful analogy proposes a falsifiable
dependency or intervention, not merely a new vocabulary.

\begin{KeyIdeas}
\begin{itemize}\setlength{\itemsep}{2pt}
\item Gene expression is a chain of physical processes (binding, transcription,
processing, translation, decay), each with its own rates and observables.
\item Cis regulation acts through a locus's own sequence; trans regulation acts
through diffusible factors. Spatial chromosome context matters.
\item A mutual-repression model has symmetric and asymmetric fixed points;
stability, not a plateau in a plot, decides which states can persist.
\item The same parameters can give different outcomes from different histories:
that is what memory means in a dynamical system.
\item Recorded lineage and inferred pseudotime are different evidence; molecule
partition at division is not DNA copying.
\item Analogies between gene regulation and computation must state their
carriers, assumptions and omitted processes.
\end{itemize}
\end{KeyIdeas}

\section{Twenty worked exercises}
\begin{enumerate}
\item Why is RNA not protein? \textbf{Answer:} translation and distinct loss/processing
dynamics separate their states; protein-only data do not identify all rates.
\item Derive one-site occupancy. \textbf{Answer:} normalize weights $1,x/K$ to get $x/(K+x)$.
\item Does Hill exponent two prove two cooperative binders? \textbf{Answer:} no; here it is phenomenological.
\item Compute the cis-A rate after halving $c_A$. \textbf{Answer:} $.25$ at $x=K=1,q_A=1$.
\item Shared trans change to $x=3$? \textbf{Answer:} both rates become $.75$ for unit cis/accessibility.
\item Why not assign an enhancer by distance alone? \textbf{Answer:} compatibility, context and perturbation evidence are distinct from linear proximity.
\item Derive $\alpha$. \textbf{Answer:} rescale concentration by $K$ and time by $1/\delta$ to obtain $v/(\delta K)$.
\item Subtract the fixed-point equations. \textbf{Answer:} $(a-b)(1-ab)=0$.
\item Asymmetric points at $\alpha=3$? \textbf{Answer:} $(3\pm\sqrt5)/2$ in opposite coordinates.
\item Local eigenvalues there? \textbf{Answer:} $-5/3,-1/3$ in dimensionless inverse time.
\item Stability of the symmetric point at $\alpha=3$? \textbf{Answer:} saddle, since its antisymmetric eigenvalue is $.19105889>0$.
\item Why does an exactly symmetric simulation not choose a branch? \textbf{Answer:} deterministic symmetry is invariant; no stochastic forcing is present.
\item Prove Euler positivity. \textbf{Answer:} for $0<h\le1$, each update is a nonnegative old-state/production sum.
\item Why reject rather than clip $h=2$? \textbf{Answer:} clipping conceals a changed, invalid numerical method.
\item Production-off versus concentration clamp? \textbf{Answer:} the former preserves initial A and lets it decay; the latter replaces its state equation.
\item Pseudotime order in the three-cell fixture? \textbf{Answer:} right, root, left; it is not ancestry.
\item Daughter covariance under independent molecule placement? \textbf{Answer:} $\operatorname{Cov}(X,N-X)=-N/4$.
\item Selection frequency after one generation? \textbf{Answer:} $.25(2)/[.25(2)+.75(1)]=.4$.
\item Frequency after six generations without mutation? \textbf{Answer:} $64/67=.955223881$.
\item Design a biology-to-model claim. \textbf{Rubric:} carrier, variables/units,
timescale, intervention, measurement, omitted processes, established computational
baseline, falsifier and independent biological evidence.
\end{enumerate}

\section{Bridge to molecular-programming evidence}
Chapter 31 returns to engineered molecular systems. A proved regulatory fixed
point is a model result; fluorescence, cell behavior and molecular circuit performance
need evidence at their own environmental and observational scope. The next chapter
asks precisely which refinement and measurement support each claimed computation.
